\documentclass[aspectratio=169]{beamer}

\usetheme{colorful-dream}

\usepackage{minted}
\setminted{fontsize=\footnotesize, breaklines=true}

\title{Basics of Python for Data Science}
\subtitle{Lecture 1: Python Basics, Environment, First Engineering Plot}
\author{Lu Pa}
\date{\today}

% --- Title Slide Logos ---
\titlegraphic{
  \vspace{0.5cm}
  \begin{center}
    \includegraphics[height=1.8cm]{logo_university_placeholder.png}
    \hspace{1cm}
    \includegraphics[height=1.8cm]{logo_faculty_placeholder.png}
  \end{center}
}

\begin{document}

% ---------------------------------------------------------
\begin{frame}
  \titlepage
\end{frame}

% ---------------------------------------------------------
\section{PART I}
\begin{frame}[plain]
  \begin{center}
    \vspace{2cm}
    {\usebeamerfont{title}\color{white}\Huge PART I}
  \end{center}
\end{frame}

% ---------------------------------------------------------
\begin{frame}{A Brief History of Python}
  \begin{itemize}
    \item Created by Guido van Rossum in 1991
    \item Designed to be readable, simple, and expressive
    \item Influenced by ABC, Modula-3, and functional languages
    \item Python 2 → 2000, Python 3 → 2008 (modern standard)
    \item Today: one of the most widely used languages in the world
  \end{itemize}
\end{frame}

% ---------------------------------------------------------
\begin{frame}{Python in Numbers}
  \begin{itemize}
    \item \textbf{\#1} most popular language in many industry surveys
    \item Over \textbf{400,000} packages on PyPI
    \item Used by:
      \begin{itemize}
        \item Google, Meta, Microsoft, Tesla, NASA, CERN
        \item Automotive, aerospace, finance, biotech, robotics
      \end{itemize}
    \item Dominant in:
      \begin{itemize}
        \item Data science
        \item Machine learning
        \item Automation and testing
        \item Scientific computing
      \end{itemize}
  \end{itemize}
\end{frame}

% ---------------------------------------------------------
\begin{frame}{Why Python? (Advantages)}
  \begin{itemize}
    \item Easy to read and write — low barrier to entry
    \item Huge ecosystem of scientific libraries
    \item Excellent community support
    \item Works across platforms (Windows, Linux, macOS)
    \item Integrates with C/C++ for performance
    \item Ideal for prototyping and production
  \end{itemize}
\end{frame}

% ---------------------------------------------------------
\begin{frame}{Limitations of Python}
  \begin{itemize}
    \item Not the fastest language (interpreted)
    \item GIL limits multi-threaded CPU-bound tasks
    \item Packaging can be confusing for beginners
    \item Not ideal for low-level embedded systems
    \item But: scientific libraries overcome most performance issues
  \end{itemize}
\end{frame}

% ---------------------------------------------------------
\begin{frame}{Python vs MATLAB / C++ / R}
  \begin{itemize}
    \item \textbf{MATLAB}: great for engineering, but expensive and closed
    \item \textbf{C++}: fast and powerful, but complex and verbose
    \item \textbf{R}: strong in statistics, weaker in general engineering
    \item \textbf{Python}: free, flexible, huge ecosystem, modern standard
  \end{itemize}
\end{frame}

% ---------------------------------------------------------
\begin{frame}{The Scientific Python Ecosystem}
  \begin{itemize}
    \item \textbf{NumPy} — fast numerical arrays
    \item \textbf{PPandas} — data tables, cleaning, grouping
    \item \textbf{Matplotlib} — plotting
    \item \textbf{SciPy} — signal processing, optimization
    \item \textbf{scikit-learn} — machine learning
    \item \textbf{Jupyter} — interactive notebooks
  \end{itemize}
\end{frame}

% ---------------------------------------------------------
\begin{frame}{Why Jupyter for Engineering?}
  \begin{itemize}
    \item Code + text + plots in one place
    \item Great for experimentation and documentation
    \item Reproducible workflows
    \item Perfect for teaching and reporting
    \item Widely used in research and industry
  \end{itemize}
\end{frame}

% ---------------------------------------------------------
\section{PART II}
\begin{frame}[plain]
  \begin{center}
    \vspace{2cm}
    {\usebeamerfont{title}\color{white}\Huge PART II}
  \end{center}
\end{frame}

% ---------------------------------------------------------
\begin{frame}[fragile]{First Imports}
\begin{minted}{python}
import numpy as np
import pandas as pd
import matplotlib.pyplot as plt
\end{minted}
\end{frame}

% ---------------------------------------------------------
\begin{frame}{The Dataset}
  \begin{itemize}
    \item Synthetic vehicle telemetry
    \item Columns:
      \begin{itemize}
        \item time\_s — time in seconds
        \item speed\_kmh — vehicle speed
        \item accel\_mps2 — acceleration
        \item engine\_rpm — engine speed
        \item fuel\_rate\_gps — fuel rate
      \end{itemize}
  \end{itemize}
\end{frame}

% ---------------------------------------------------------
\begin{frame}[fragile]{Loading the Data}
\begin{minted}{python}
df = pd.read_csv("../datasets/vehicle_speeds_trackA.csv")
df.head()
\end{minted}
\end{frame}

% ---------------------------------------------------------
\begin{frame}[fragile]{First Engineering Plot}
\begin{minted}{python}
plt.figure(figsize=(10,4))
plt.plot(df["time_s"], df["speed_kmh"])
plt.xlabel("Time [s]")
plt.ylabel("Speed [km/h]")
plt.grid(True)
plt.title("Vehicle Speed Over Time (Track A)")
plt.show()
\end{minted}
\end{frame}

% ---------------------------------------------------------
\begin{frame}{Practice Tasks}
  \begin{itemize}
    \item Convert speed from km/h to m/s
    \item Compute min, max, and average speed
    \item Plot RPM vs time
    \item Highlight maximum speed on the plot
    \item Filter data for speeds above 50 km/h
    \item Load Track B and compare profiles
  \end{itemize}
\end{frame}

% ---------------------------------------------------------
\end{document}
